\begin{exercisechunk}
\item \textbf{Which correlations determine the gradient?}
\label[exercise]{ex:l8-gradient-features}
\exerciseprofile{2}{2}{3}{\exproof\quad\excalculation\quad\exinterpretation}{%
Identify the features that govern a local gradient step and quantify the
sampling error in estimating their correlations.}
Use the ReLU network, fixed parity $Y=\chi_S(X)$, and objective $F$ from
\cref{sec:recap}. Fix $\theta$ such that $w_r^\top x+b_r\ne0$ for all
$r\in[m]$ and $x\in\{-1,1\}^d$. Define the vector of
\emph{parameter derivative features}
\[
\Phi_\theta(x)=\nabla_\theta N_\theta(x)\in\R^p.
\]
\begin{enumerate}[(a),beginpenalty=10000]
\item Prove
\[
\max_{\|v\|\le1}\lim_{t\downarrow0}
\frac{F(\theta)-F(\theta+tv)}{t}
=\|\E[Y\Phi_\theta(X)]\|.
\]
Give a maximizing $v$, including the case when the expectation is zero.
Explain the meaning of an upper bound $\|\E[Y\Phi_\theta(X)]\|\le\varepsilon$
for first-order improvement in a unit parameter direction.

\item Write the coordinates of $\Phi_\theta(x)$ corresponding to
$u_r$, $w_{rj}$, and $b_r$. Express their correlations with $Y$ using
$q_B(w,b)=\E[\chi_B(X)\I\{w^\top X+b>0\}]$:
\[
\begin{aligned}
\E[Y\Phi_{\theta,u_r}(X)]
 &=\sum_{j=1}^d w_{rj}q_{S\mathbin\triangle\{j\}}(w_r,b_r)
   +b_rq_S(w_r,b_r),\\
\E[Y\Phi_{\theta,w_{rj}}(X)]
 &=u_rq_{S\mathbin\triangle\{j\}}(w_r,b_r),\\
\E[Y\Phi_{\theta,b_r}(X)]&=u_rq_S(w_r,b_r).
\end{aligned}
\]
Identify which coordinates are the hidden-unit outputs and which involve
the activation indicators.

\item Fix $w\in\R^d$ such that $w^\top x\ne0$ for every cube point $x$,
and fix $j\in[d]$. Let $X^{(1)},\ldots,X^{(n)}$ be independent uniform
cube points and set
\[
q=q_{S\mathbin\triangle\{j\}}(w,0),
\qquad
\widehat q_n=\frac1n\sum_{i=1}^n
\chi_S(X^{(i)})X_j^{(i)}\I\{w^\top X^{(i)}>0\}.
\]
Prove
\[
\E[(\widehat q_n-q)^2]=\frac{1/2-q^2}{n}.
\]
If $0<|q|\le e^{-ck}\le1/2$, where $c>0$, prove that the requirement
$\sqrt{\E[(\widehat q_n-q)^2]}\le |q|$ implies
$n\ge e^{2ck}/4$.
\end{enumerate}
\begin{hints}
For (a), differentiate along the line $\theta+tv$.
For (b), use $[z]_+=z\I\{z>0\}$ and
$\chi_S(x)x_j=\chi_{S\mathbin\triangle\{j\}}(x)$.
For (c), pair $x$ with $-x$ when calculating the second moment of one summand.
\end{hints}
\end{exercisechunk}
